% !TEX root = ../main.tex
\section{Introduction}
\label{sec:intro}

This section checks the basic text features. A paragraph is written as ordinary continuous text, and a single line break in the source does not change anything in the PDF,
because LaTeX only starts a new paragraph at a blank line.

A second paragraph should appear after a blank space and without indentation, which shows that the parskip package is active.
The first page should have no header, a rule above the footer, and the three author emails at the bottom left.
Clicking a symbol next to an author name should jump to the corresponding email, and clicking an email should open your mail program.
From the second page on, the header should show the assignment on the left and the course on the right, with a thick rule below.

The bibliography works if the following citations appear as numbers in brackets and not as question marks.
Hawking showed that black holes emit thermal radiation~\cite{hawking1975particle}, and the classical theory of the oscillator used below can be found in any electrodynamics or mechanics text~\cite{jackson1999classical,goldstein2002classical}.
The citation links should be dark red.

Cross-references work if the following sentence contains working blue links with the right names and numbers.
The model is derived in \cref{sec:theory}, its solution is \cref{eq:solution}, and the results are in \cref{fig:oscillator,tab:parameters}.

Notes between co-authors appear as colored boxes inside the text:
\todo{This is an example note. To hide every note, add the disable option to todonotes at the top of main.tex.}

External links should be red, for example the \href{https://tug.org/texlive/}{TeX Live home page}.
